Recent coding agents have made tool-using, stateful execution prominent.
The same execution style is increasingly used for knowledge work that ends in long-form documents.
Yet long-form writing systems usually advance an outline, task graph, or fixed stage sequence rather than deciding which writing process is needed next.
Revisiting cognitive process theory as a framework for process-level control in AI writing agents, we propose \condAgenticCogWriter, which maintains a shared draft, goals, and process history while a Monitor agent chooses and delegates planning, drafting, or reviewing.
We evaluate our \condAgenticCogWriter\ across three long-form writing benchmarks using pointwise rubric scoring and pairwise same-prompt preference judgments.
\condAgenticCogWriter\ achieves competitive pointwise quality and is preferred to one-shot generation, a fixed planning--drafting--revision workflow, and adaptive task planning in pairwise evaluation.
Ablations indicate that neither explicit goals nor adaptive process order alone explains \condAgenticCogWriter's overall advantage.
Process traces show that most runs follow one planning--drafting--reviewing cycle, with departures from that path used selectively rather than continually.
Together, the results support process-level control as a promising design for long-form writing agents.
